\egasection{§15\quad 前述若干结果的公理化}

% SOURCE: unpacked/dot_batches/EGA第五卷_第15节公理化导言_续译工程/EGA第五卷_第15节公理化导言_校订续译.tex
我总的看法是，第2至第8节的结果大多在比射影空间中的超平面族（或给定次数的超曲面族）更一般的条件下仍然成立。采用公理化的观点似乎是合适的。\cmt{V-H-L048}{底本首句有“tht”等拼写及句法缺陷，正文按其可辨意义译出。“大多”保留“mostly”的限定，不把第2至第8节的全部结果无条件推广。}
我目前还不十分确定，能否在这个意义上推广Bertini--Zariski（从而推广第4节和第6节的结果）；我已经写信给这方面的权威（译者：有能力的人），即Serre和Zariski，询问他们是否知道这样的推广。不过，我的确有这样的印象：上述那类简单的微分性质的假设，应当足以推出Bertini--Zariski。如果这些专家（有能力的人）不能给我们满意的答复，我们就应当设法用自己的办法把这件事弄清楚。\cmt{V-H-L049}{底本的“(competent people Tr)”以及“(experts) competent people”作为编辑痕迹译出。“我”指原作者；写信及请教的叙述不是本译者已实施的行动。原文“not quite sure”及“impression”均保留，不据此宣称Bertini--Zariski推广已建立。}
我们从一个由有限表现态射组成的交换图式出发：\cmt{V-H-L050}{底本图式上方两个对象都印为$P$。紧接第60页说明其中一个为射影纤维空间、另一个为关联预概形，并讨论它到$P\times_S G$的态射。故正文保留左上$P$，将右上关联预概形记作$\mathcal P$，以免误读为同一对象的自映射。这里仍只是交换图式，没有把它擅自加强成笛卡尔图式。}
\[
\begin{tikzcd}
P \arrow[d] & \mathcal P \arrow[l] \arrow[d] \\
S & G \arrow[l]
\end{tikzcd}
\]

% SOURCE: unpacked/dot_batches/EGA第五卷_参数族与基变换_续译工程/EGA第五卷_参数族与基变换_校订续译.tex
（在主要应用中，$P$是射影纤维空间，$G$是由它导出的Grassmann空间〔“grassmannienne”是形容词！〕\footnote{底本脚注68：可能是Grassmann纤维空间。[原英译者注]}，$\mathcal P$是关联预概形。在最重要的情形中，相应态射
\[
\mathcal P\longrightarrow P\times_S G
\]
应当是闭浸入；我们把$G$看作$P$在$S$上的各纤维中的一族闭子预概形的参数概形。更确切地说，若$\xi\in G$，则$\mathcal P_\xi$是$P_s\times_{k(s)}k(\xi)$的闭子预概形，其中$s$是$\xi$在$S$中的像。此外，对于这里的大多数陈述，我们无疑取$S=\operatorname{Spec}(k)$。）\cmt{V-H-L051}{底本写“$P\to P_S\times G$”，正文依上一页交换图式规范为$\mathcal P\to P\times_S G$。底本“$s$ is the point of $S$ over $\xi$”把方向倒置，按结构态射$G\to S$改为$s$是$\xi$的像。$P_s,k(\xi)$明确写为剩余域扩张后的纤维，不把$\mathcal P_\xi$直接当作$P_s$上的闭子概形。}
在一般情形中，仍可把$G$看作$P$在$S$上的各纤维之上的一族预概形的参数概形，$\xi$对应于$P_s\times_{k(s)}k(\xi)$之上的$\mathcal P_\xi$。当然，不必把$\xi$取为$G$的一个（绝对）点，也可以取它为一个以$S$预概形$T$为值的点。于是得到
\[
\mathcal P_\xi\longrightarrow P_T
\]
这一$T$态射；在上述特殊情形中，它是闭浸入。\cmt{V-H-L052}{底本此处目标写作$B_T$，而本段及交换图式没有引入$B$。由$\xi:T\to G$拉回$\mathcal P\to P\times_S G$，其目标是$P_T$，故正文订正为$P_T$并保留原写法于此。这里$T$值点不要求$T$为一个域的谱。}


若$f:X\to P$是态射，令
\[
\mathcal X=X\times_P\mathcal P,
\]
就得到与前面的方形同类型的图式：\cmt{V-H-L053}{底本写“we put $X=X_P\times P$”，图式上方两个对象也都印作$X$。依沿$f$作基变换的构造，正文用$\mathcal X$表示新的纤维积，并规范下标位置；原来的$X$保持不变。此处所画方形仅断言交换，不另加笛卡尔性。}
\[
\begin{tikzcd}
X \arrow[d] & \mathcal X \arrow[l] \arrow[d] \\
S & G \arrow[l]
\end{tikzcd}
\]
因此，第2至第8节研究的所有问题，在我们刚刚陈述的一般背景下仍有意义；并且有充分理由去明确那些保证上述各节结论成立的公理化条件。\footnote{底本脚注69：“Fr dégager + unravel, make explicit ?? [Tr]”，即法语dégager，疏解、明确表述？？[原英译者注]。正文中的“明确”据此恢复底本“there is a good reason to”之后缺失的动词含义，保留词义选择的问号。}\cmt{V-H-L054}{原句有“It is therefore evident”这一概括用语；译文写“因此”承接前面的构造，不把该句当作全部推广结论的证明。作者在这里提出的是寻找充分公理条件的任务，条件本身尚待下文列出。}

% SOURCE: unpacked/dot_batches/EGA第五卷_平坦性假设与非空性公理_续译工程/EGA第五卷_平坦性假设与非空性公理_校订续译.tex
我们将假设$P$和$G$在$S$上平坦，且$G$的纤维几何不可约、维数为$N$（以便能够考虑泛点！）；假设态射$\mathcal P\to P$光滑，其纤维几何不可约、维数为$N-m$。于是态射$\mathcal X\to X$也具有同样的性质。\cmt{V-H-L055}{沿用$\mathcal X=X\times_P\mathcal P$的记号；底本将这两个态射分别印作$P\to P$和$X\to X$。正文按其基变换关系区别源与目标，不改动光滑性、几何不可约性及维数假设。}
所有提到的性质〔字迹不清，且内容较长〕\footnote{底本脚注70：后面的字迹不清。}在$S$上的基变换下都是稳定的（保持不变），特别可以应用于纤维。\cmt{V-H-L056}{底本括注为“[(illegible and long)]”。缺失部分不补写；原文对基变换稳定性的概括照译，不把未能辨出的内容视作已经逐项核验。}

首先设$S=\operatorname{Spec}(k)$。令$Z$为$X$的一个$d$维闭子集；于是它在$\mathcal X$中的逆像$\mathcal Z$是一个维数为
\[
d+(N-m)=N+d-m
\]
的闭子集。\cmt{V-H-L057}{底本把闭子集及其逆像都印为$Z$，正文以$\mathcal Z$区别逆像。原文直接使用这一维数等式；若允许$\mathcal P\to P$有空纤维，还须注意逆像非空及所涉及分支是否遇到其像，不能仅凭“光滑”无条件宣称任意闭集的逆像具有该维数。这里保留原论证并标明适用性待核。}
若$d<m$，则$\mathcal Z$的维数小于$N$，所以$\mathcal Z\to G$不可能占优。因此，若$\eta$是$G$的泛点，就有$\mathcal Z_\eta=\varnothing$。


事实上，这个推理还表明（将$Z$换成$\overline{f(Z)}$）：若$\dim f(Z)<m$，则$\mathcal Z_\eta=\varnothing$。\cmt{V-H-L058}{底本括号内印有$f(Z)$的上划线，正文保留闭包。紧接的条件原为“$\dim f(X)<m$”；这本身是更强的充分条件，但与“将$Z$换成$\overline{f(Z)}$”及下句所讨论的$\dim f(Z)$不对应，故按本段论证对象改为$\dim f(Z)<m$，并在此保留原写法。仍沿用底本对像维数的记法，不引入新的证明。}
我们希望对图式$(D)$提出一个条件，以保证当$\dim f(Z)\geq m$时，$\mathcal Z_\eta\ne\varnothing$。看来，在这里应当把下述要求作为一条基本公理（在2.2的情形中，它会由一个相当特殊的整体论证得出）：对$P$的每个维数为$m$的闭不可约子集$Z$，都有\cmt{V-H-L059}{底本末句实际印作“$Z_\eta=\phi$”。但紧邻前句明确要求$\dim f(Z)\geq m$时的非空性，而第61页又据此讨论占优性；故此处等号为与段意相反的误写，正文订为$\ne\varnothing$。本句$Z$改取为$P$中的闭集，$\mathcal Z$相应指其在$\mathcal P$中的逆像。这里是作者建议加入的公理，不是已由前述光滑、平坦假设推出的定理；从维数恰为$m$的闭集推广到其他情形的细节尚未在本段译校补证。}
\[
\mathcal Z_\eta\ne\varnothing.
\]

% SOURCE: unpacked/dot_batches/EGA第五卷_泛截面光滑性条件_续译工程/EGA第五卷_泛截面光滑性条件_校订续译.tex
再取$X$的闭子集$Z$，使得$\dim f(Z)\geq m$。我们看到$\mathcal Z\to G$占优，因而$\mathcal Z_\eta$的维数为\cmt{V-H-L060}{底本写“$Z_\eta\to G$ is dominant”，继而说“$Z$ is of dimension equal to $\dim Z-\dim G=\dim Z-m$”。按上页$\dim\mathcal Z=N+\dim Z-m$及$\dim G=N$，占优的应是总空间$\mathcal Z\to G$，维数等式左边应指泛纤维$\mathcal Z_\eta$。正文据此区别总空间、原闭集和泛纤维。闭集可能有多个分支，哪些分支占优以及前文逆像维数等式的适用性仍须核对，不据此宣称任意闭集的全部分支均满足同一纯维数结论。}
\[
\dim\mathcal Z-\dim G=\dim Z-m.
\]
这些性质使我们能够在当前背景下展开与2.1和2.11相应的结果。还有一个关于〔字迹不清〕的条件（保证2.12成立，即假设$f:X\to P$非分歧时，若$X$光滑，则相应截面也光滑）。\cmt{V-H-L061}{底本括注写“if $X$ is smooth then $X$ is also such”，两个对象没有区别。结合本段最终结论$\mathcal X_\eta$光滑，这里译作“相应截面”；不把它解释成整个族$\mathcal X\to G$处处光滑，亦不补出不可辨处的条件。}

现在假设$P$在$k$上光滑，$\mathcal P\to P\times_kG$拟有限，并满足以下条件（这里假设$k$代数闭）。对于每个$x\in P(k)$和切空间$T_x(P)$的每个维数$n\geq m$的向量子空间$V$，考虑由如下$\xi\in G(k)$组成的集合$E(x,V)$：$\mathcal P_\xi$中存在位于$x$之上的点$z$，它不满足以下整组条件：\cmt{V-H-L062}{底本先说“$P$ has a point over $x$”，下文才出现$z$和$P_\xi$；切空间源又印作$T_z(P)$。正文依所指定的纤维态射统一为$z\in\mathcal P_\xi$及$T_z(\mathcal P_\xi)$。“不满足以下整组条件”表示三项不全成立，不是要求三项全部失败。}
\begin{itemize}
\item $\mathcal P_\xi$在$z$处光滑；
\item $\mathcal P_\xi\to P$在$z$处诱导的切映射
\[
T_z(\mathcal P_\xi)\longrightarrow T_x(P)
\]
是单射（即$\mathcal P_\xi\to P$在$z$处非分歧）；
\item 该切映射的像与$V$“横截”，即二者之和为$T_x(P)$。
\end{itemize}



所要求的条件是：$E(x,V)$的维数$\leq N-n-1$（我们知道它是$G$中一个良定义的可构造集在$G(k)$上的迹，???）。\cmt{V-H-L063}{底本的三个问号原样保留。正文没有据此宣布可构造性或坏参数集的维数估计已经证明；该维数上界是本段要求加入的条件。}
借助〔Moyennant〕这个条件，应用雅可比判据并作维数计算，可知由下列点构成的闭子集$E$的维数不超过
\[
n+(N-n-1)=N-1:
\]
这些点是$\mathcal X$中使$\mathcal X\to G$不光滑的点，或其在$\mathcal P$中的像使$\mathcal P\to G$不光滑的点，或其像使原文所写“$P\to P$”分歧的点（这里$X$处处光滑、维数为$n$）。\cmt{V-H-L064}{底本将$\mathcal X$与$X$以及$\mathcal P$与$P$均印作同一字母，并将诱导映射$\mathcal X\to\mathcal P$仍记作$f$，正文明确了像所在的空间。但最后“$P\to P$ is ramified”究竟指总空间态射还是各纤维态射不能直接确定：前面假设的是$\mathcal P\to P$光滑，而本段切映射条件涉及$\mathcal P_\xi\to P$非分歧。故此处保留原写法并标歧义，不将它默改为一个已验证的坏点集定义。该集合的闭性和维数估计尚待连同缺文核对。}
因此$\dim E<N=\dim G$，从而$E_\eta=\varnothing$；特别地，$\mathcal X_\eta$在$k(\eta)$上光滑，第5节的论述在当前背景下也有效。\cmt{V-H-L065}{底本“$N=\dim$”缺少对象，按前文$\dim G=N$补为$N=\dim G$。这段是在所设条件下的原论证概要，未展开雅可比判据和维数计数；上述歧义不因翻译完成而消除。}

% SOURCE: unpacked/dot_batches/EGA第五卷_第15节推广计划结语_续译工程/EGA第五卷_第15节推广计划结语_校订续译.tex
第4节中从泛截面到一般截面的过渡，以及第5节的论述，在当前背景下也有效（不过，到这里，它们已成为同义反复，或只是对第8、9、12节的重新表述；我们对于以这种形式陈述它们仍有所犹豫）。\cmt{V-H-L066}{底本有“evidently valid”，此处按原句结论译为“也有效”，不将其当作本译者已补全的一般化证明。“paragraphs 8, 9, 12”沿用本稿编号体系译为第8、9、12节。}
同样，若$k$代数闭，7.1的论述在所有情形中都有效；若再假设$G$在$k$上有理，那么即使只要求$k$无限也成立。特殊情形7.2、7.3$'$也是如此；至于〔quant，???〕7.4的结果，则是针对超平面截面情形的特殊应用。如我已经说过的，第3节和第6节仍须暂缓，等待Zariski定理的推广。\cmt{V-H-L067}{底本原句为“and the special cases 7.2, 7.3' quant (???) the result 7.4”。正文保留7.3后的撇号和“quant”旁的问号，仅按句意以“至于”作暂定连接，不把编号改成另一个确定条目。“rational over $k$”译为“在$k$上有理”，不是指$G$仅有一个$k$有理点；括号“suspended pending”明确保留为等待推广的状态。}

还须推广第8节的结果（第12节曾重新考虑过），它们由此呈现出更令人满意的面貌。我建议你先在这个背景下表述这些结果，并尝试沿这条路尽可能推进。我觉得，我们至少应当能够重新得到那些不是8.7(c)直接推论的结果（甚至可以尝试提炼出一些公理化条件，使你能够完成8.7的一个变式）。\cmt{V-H-L068}{底本此句实际为“recover at least that is not a direct consequence of 8.7 c)”，主语或先行词不完整，正文补出最低限度的“那些……结果”以成句；其中“not”照译为“不是”，没有擅自删去否定词。底本括号在下一句前未闭合，正文按语义在“变式”后补右括号。原文究竟意欲恢复哪些结果及否定词是否可靠仍待核，不据此修复8.7(c)的历史缺文。}
我这里只提出这些建议；不过，若你遇到特殊困难，我愿意再回到这些问题上，给出更多细节。\cmt{V-H-L069}{这里的“我”和“你”是原稿作者与收信人之间的叙述，保留建议、印象及愿意进一步说明的语气，不把它改写成本译者已经执行的研究或承诺。}
